% ===================================================================
%  TAFEL Nr. 16 — Das Kaufhaus. --- Der Basar. --- Das Spielzeug.
%  Traduction 2026 d'apres texto/io, controlee sur texto/fr et
%  texto/en. Consigne : voir texto/de/10-tafel-01.tex.
%
%  Le tableau d'astronomie (bloc 15-1) porte des renvois a LETTRES,
%  (a) a (f), graves sur le tableau lui-meme : ils se rangent sous le
%  numero 85, comme le dit gravuri/literi.json. L'ido en compose un
%  avec une seule parenthese, « d) » ; l'allemand met les deux.
%
%  CETTE PAGE CORRIGE UNE CONCLUSION QUE LA SERIE AVAIT TIREE TROP
%  VITE, ET C'EST L'ALLEMAND QUI PEUT LE FAIRE PARCE QU'IL EST LA
%  SOURCE.
%
%  Le tableau 16 romanche avait trouve que la menagerie de carton de
%  l'alinea 3 etait nommee par le JOUET, arrive de Nuremberg avec son
%  etiquette imprimee. Le tableau 16 estonien avait ensuite trouve que
%  l'estonien, lui, « composait tout » — ninasarvik le corne-au-nez,
%  jõehobu le cheval-de-riviere, kilpkonn la grenouille-a-bouclier,
%  nahkhiir la souris-de-cuir — et l'on avait conclu qu'il avait
%  REPEINT l'etiquette.
%
%  IL NE L'A PAS REPEINTE. IL L'A TRADUITE. Voici l'allemand :
%     Nashorn      nez-corne        estonien ninasarvik  nez-corne
%     Flusspferd   cheval-de-fleuve           jõehobu    cheval-de-riviere
%     Schildkröte  crapaud-a-bouclier         kilpkonn   grenouille-a-bouclier
%     Fledermaus   souris-volante             nahkhiir   souris-de-cuir
%  Morpheme pour morpheme, dans le meme ordre, et l'ordre n'est pas
%  celui de l'estonien ordinaire mais celui de l'allemand. Le critere
%  du tableau 12 tranche : « nez-corne » est une image evidente, on
%  peut la trouver seul ; « crapaud-a-bouclier » ne l'est pas, et
%  deux langues sans rapport ne tombent pas dessus par hasard. C'EST
%  UN CALQUE, ET IL VIENT DE NUREMBERG COMME LE JOUET.
%
%  La serie s'etait donc trompee deux fois de suite sur le meme
%  alinea, et pour la meme raison : elle ne connaissait pas encore la
%  langue de l'etiquette. On l'inscrit ici plutot que de reecrire les
%  deux en-tetes precedents, qui restent le compte rendu exact de ce
%  qu'on pouvait voir a ce moment-la.
%
%  ET C'EST LA FIN DE LA COLONNE ALLEMANDE. Elle aura trouve cinq
%  facons pour un mot de passer d'une langue a une autre, et chacune
%  a sa page :
%     1. le mot arrive avec la chose             (Gletscher, t. 9)
%     2. le mot est decrete par quelqu'un        (Turnen, t. 2 ; la
%        marine, t. 10 ; la Feuerwehr, t. 11)
%     3. la forme est calquee, non le mot        (Bahnhof, t. 12 ;
%        et ici la menagerie)
%     4. l'institution garde la langue du pays qui l'a batie (t. 11)
%     5. on garde le mot etranger parce qu'etre etranger est ce
%        qu'il sert a dire                       (l'hotel, t. 13)
%  Et une regle qui les traverse toutes : ON PEUT DECIDER DU NOM D'UNE
%  CHOSE QUI ARRIVE ; ON NE PEUT PAS REPRENDRE LE NOM D'UNE CHOSE QUI
%  EST LA (t. 3).
%
%  Un dernier compose, et il est de la meme fabrique que la menagerie :
%  la boussole de l'alinea 15 nomme ses quatre points Norden,
%  Süden, Osten, Westen — les mots
%  germaniques que toute l'Europe a empruntes aux marins du nord, et
%  que le tableau 16 estonien avait justement REFUSES, preferant son
%  põhi, lõuna, ida, lääs. La derniere page de la premiere serie et la
%  derniere page de celle-ci se repondent sur le meme objet : l'un a
%  pris la rose des vents des marins, l'autre l'a nommee par le soleil
%  et par l'heure du repas.
% ===================================================================

%%K t16-apar-1 apar
\VUsaut{4.0mm}
\VUtitre{15.0pt}{60}{TAFEL Nr. 16}
\VUsaut{2.0mm}
\VUpk{13.2pt}{40}{Das Kaufhaus. --- Der Basar.}
\VUpk{13.2pt}{40}{Das Spielzeug.}
\VUsaut{3.4mm}

%%K t16-01-1 p
1. --- Das neue Jahr steht vor der Tür. Die großen Läden haben ihre Auslagen mit
\VUgras{Spielzeug} \textsuperscript{(1)} und Neujahrsgeschenken hergerichtet.

%%K t16-01-2 p
Ich habe meinen Großvater gebeten, mich in den Basar mitzunehmen, um diese schöne
Auslage zu sehen, und er war einverstanden.

%%K t16-02-1 p
2. --- Zuerst blieben wir vor ein paar schönen Waffentrophäen stehen, mit einem
\VUgras{Gewehr}, einer \VUgras{Dienstmütze}, einem \VUgras{Säbel}, einem
\VUgras{Degengehänge} und einem Paar \VUgras{Epauletten}, oder auch mit einem
\VUgras{Helm}, einem \VUgras{Brustpanzer} \textsuperscript{(3)} und einer
\VUgras{Pistole} \textsuperscript{(4)}. Das alles hat mich weit mehr interessiert als
die \VUgras{Laterna magica} \textsuperscript{(2)}.

%%K t16-tit-1 sub
\VUsaut{3.0mm}
\VUpk{10.2pt}{40}{Schöne Sachen.}
\VUsaut{2.6mm}

%%K t16-03-1 p
3. --- In derselben Abteilung stand ein \VUgras{Posten} \textsuperscript{(5)} in
seinem \VUgras{Schilderhaus}; er schien über eine ganze Menagerie aus Pappe zu
wachen. Was für eine Menge Tiere! Ein \VUgras{Papagei} \textsuperscript{(6)} auf
seiner Stange, ein \VUgras{Nashorn} \textsuperscript{(7)} mit einem Horn auf der
Nase, ein \VUgras{Rentier} \textsuperscript{(8)}, ein \VUgras{Strauß}
\textsuperscript{(9)}, ein \VUgras{Affe} \textsuperscript{(10)}, der eine Nuss
knackt, ein \VUgras{Fuchs} \textsuperscript{(11)}, der eine Henne davonträgt, ein
\VUgras{Krokodil} \textsuperscript{(12)}, das seinen riesigen Rachen aufreißt, ein
\VUgras{Kamel} \textsuperscript{(13)}, ein eleganter \VUgras{Windhund}
\textsuperscript{(14)}, ein \VUgras{Panther} \textsuperscript{(15)} und dann zwei
Tiere, die ich nicht kannte und die, wie man mir sagte, ein \VUgras{Wiesel}
\textsuperscript{(16)} und eine \VUgras{Hyäne} \textsuperscript{(17)} sind. Es gab
auch einen \VUgras{Leoparden} \textsuperscript{(18)} und einen \VUgras{Wolf}
\textsuperscript{(21)}; gleich daneben lagen niedliche kleine \VUgras{Ratten}
\textsuperscript{(19)} und winzige \VUgras{Mäuse} \textsuperscript{(20)} aus Gummi.
Zuletzt vervollständigten ein \VUgras{Flusspferd} \textsuperscript{(22)} und ein
\VUgras{Dromedar} \textsuperscript{(23)} mit seinem Höcker die Sammlung. Ich wäre
stundenlang dort geblieben, wenn nicht ganz in der Nähe ein prächtiges Lager voller
\VUgras{Zinnsoldaten} \textsuperscript{(29)} mir ins Auge gefallen wäre.

%%K t16-tit-2 sub
\VUsaut{4.0mm}
\VUpk{10.2pt}{40}{Soldaten und Krieg.}
\VUsaut{2.8mm}

%%K t16-04-1 p
4. --- Mein Großvater, der ein \VUgras{Kriegsversehrter} \textsuperscript{(24)} ist
und nach einer \VUgras{Verwundung}, die ihm die \VUgras{Militärmedaille}
eingebracht hat, aus der Armee ausscheiden musste, erzählt mir gern von den
\VUgras{Feldzügen}, an denen er teilgenommen hat. Er erklärte mir mit Vergnügen die
verschiedenen Bewegungen des kleinen Heeres im Kleinen. Eine Gruppe wirklicher
Soldaten, die den Basar besuchten --- ein \VUgras{Pionier} \textsuperscript{(25)}, ein
\VUgras{Gebirgsjäger} \textsuperscript{(26)} mit seiner \VUgras{Baskenmütze}, ein
\VUgras{Feldwebel} \textsuperscript{(27)} der Infanterie und ein
\VUgras{Wachtmeister} \textsuperscript{(28)} der Artillerie mit einer goldenen
\VUgras{Litze} am Ärmel --- blieb neben uns stehen, um zuzuhören.

%%K t16-05-1 p
5. --- Mein Großvater, sehr stolz auf ein so vornehmes Publikum, entwarf aus dem
Stegreif einen ganzen Schlachtplan.

%%K t16-05-2 p
Die befestigte Stadt mit ihren \VUgras{Wällen} und \VUgras{Gräben} wurde vom
\VUgras{feindlichen Heer} belagert, das sich nach einem langen
\VUgras{Bombardement} zum \VUgras{Sturm} bereitmachte. Aber die
\VUgras{Belagerten}, vom Hunger dazu getrieben, hatten einen \VUgras{Ausfall} gegen
das \VUgras{Lager} der \VUgras{Belagerer} versucht. Nachdem die
\VUgras{siegreichen} Belagerer den \VUgras{Angriff} in einem heftigen
\VUgras{Kampf} um die \VUgras{Zelte} abgewehrt hatten, schickten sie ihre
\VUgras{Schwadronen} zur \VUgras{Verfolgung} der \VUgras{geschlagenen} Angreifer.
Diese \VUgras{zogen sich zurück} und ließen das \VUgras{Schlachtfeld} mit
\VUgras{Toten} und \VUgras{Verwundeten} bedeckt zurück. \VUgras{Krankenträger}
brachten die Verwundeten ins \VUgras{Feldlazarett}, damit der \VUgras{Chirurg} sie
versorge.

%%K t16-05-3 p
Trotz des heldenhaften Widerstands einiger \VUgras{Bataillone}, die das
\VUgras{Karree} gebildet hatten, war die \VUgras{Niederlage} vollständig; der Rest
des Heeres \VUgras{floh} und ließ viele \VUgras{Gefangene} zurück. Der Feind war im
Begriff, seinen \VUgras{Sieg} mit der Einnahme der Stadt zu krönen. Vielleicht wäre
das das Ende des Feldzugs, und nach einem \VUgras{Waffenstillstand} könnte ein
\VUgras{Friedensvertrag} geschlossen werden.

%%K t16-tit-3 sub
\VUsaut{4.4mm}
\VUpk{10.2pt}{40}{Neujahrsgeschenke.}
\VUsaut{2.8mm}

%%K t16-06-1 p
6. --- Die jungen Soldaten, die über die farbige Erzählung des Veteranen lachten,
stellten ihm noch ein paar Fragen. Ich ging unterdessen um den Ladentisch herum, um
mir eine schöne \VUgras{Armbrust} \textsuperscript{(32)} und einen \VUgras{Bogen}
\textsuperscript{(33)} mit seinem \VUgras{Pfeil} \textsuperscript{(31)} anzusehen.
Dort fand ich noch eine Sammlung Tiere: zwei \VUgras{Singvögel}
\textsuperscript{(34)}, eine \VUgras{Nachtigall} und eine \VUgras{Grasmücke} aus
Uhrwerk, die aus vollem Hals sangen, und eine Reihe seltsamer Tiere --- eine
\VUgras{Fledermaus} \textsuperscript{(35)}, eine \VUgras{Boa} \textsuperscript{(36)},
eine \VUgras{Schildkröte} \textsuperscript{(37)}, ein \VUgras{Seehund}
\textsuperscript{(38)}, ein \VUgras{Wal} \textsuperscript{(39)} und ein
\VUgras{Hai} \textsuperscript{(40)} aus Goldschlägerhaut.

%%K t16-07-1 p
7. --- Ich blieb nicht lange bei ihnen, weil ich das erblickt hatte, wovon ich
träumte: ein prächtiges \VUgras{Puppentheater} \textsuperscript{(41)} mit schönen
\VUgras{Kulissen} und einem entzückenden roten \VUgras{Vorhang}. Auf der
\VUgras{Bühne} schienen zwei Schauspieler eine \VUgras{Rolle} in einem sehr
spannenden \VUgras{Stück} zu spielen, denn die kleinen \VUgras{Zuschauer} aus
Pappe in den Logen oder auf dem \VUgras{Parkett} und im
\VUgras{Stehplatz} hinter dem \VUgras{Dirigenten} klatschten kräftig Beifall.

%%K t16-07-2 p
Leider war dieses Theater sehr teuer.

%%K t16-08-1 p
8. --- Die Wahl unter den Spielsachen war ohnehin schwer, weil die Schaufenster mit
Spielzeug jeder Art vollgestapelt waren: \VUgras{Harlekin} \textsuperscript{(42)},
\VUgras{Kasperle} \textsuperscript{(43)}, \VUgras{Pierrot} \textsuperscript{(44)},
\VUgras{Eulen} \textsuperscript{(45)}, die mit den Flügeln schlagen, pfeifende
\VUgras{Amseln} \textsuperscript{(46)}, bellende \VUgras{Pudel}, ein
\VUgras{Grammofon} \textsuperscript{(47)} und \VUgras{Geduldspiele}. Eines dieser
Spielzeuge hat mir besonders gefallen: es zeigte eine \VUgras{Seeschlacht}
\textsuperscript{(48)}, in der ein \VUgras{Schiff} von \VUgras{Seeräubern}
angegriffen wird, vor einer mit \VUgras{Kokospalmen} bepflanzten Küste.

%%K t16-noto-1 noto
\VUnotes{143.2mm}{%
\textit{(*) Siehe Tafel Nr. 6.}%
}

%%K t16-09-1 p
9. --- Ich konnte all diese Wunder ganz in Ruhe ansehen, weil nur wenige Leute im
Laden waren --- kaum eine Handvoll Müßiggänger, ein \VUgras{Dragoner}
\textsuperscript{(49)} und ein \VUgras{Kassierer} \textsuperscript{(50)}, der an der
Haupt\VUgras{kasse} \textsuperscript{(51)} einen \VUgras{Wechsel} vorlegte.

%%K t16-10-1 p
10. --- Ich fragte meinen Großvater, wozu die Bücher auf den Regalen hinter der
\VUgras{Kassiererin} da seien; er sagte mir, es seien die \VUgras{Geschäftsbücher}.

%%K t16-tit-4 sub
\VUsaut{3.6mm}
\VUpk{10.2pt}{40}{Ein Rundblick durch den Laden.}
\VUsaut{2.8mm}

%%K t16-11-1 p
11. --- Als wir die Spielzeugabteilung verließen, gingen wir zur Sattlerei, um einen
Blick auf die \VUgras{Sättel} \textsuperscript{(53)}, \VUgras{Steigbügel}
\textsuperscript{(54)}, \VUgras{Zäume} \textsuperscript{(55)}, \VUgras{Gebisse}
\textsuperscript{(56)} und \VUgras{Reitgerten} zu werfen. Während der
\VUgras{Abteilungsleiter} \textsuperscript{(57)} meinem Großvater Auskunft gab, sah
ich mir die Auslage mit \VUgras{Porzellan} und \VUgras{Glaswaren}
\textsuperscript{(58)} an, wo es schöne \VUgras{Salatschüsseln} und zierliche
\VUgras{Teekannen} \textsuperscript{(59)} gab.

%%K t16-12-1 p
12. --- Um in den ersten Stock zu kommen, gingen wir hinter dem Schaufenster mit den
\VUgras{Korsetts} \textsuperscript{(60)} vorbei, neben der Strumpfwarenabteilung, wo
sehr feine \VUgras{Unterhosen} \textsuperscript{(63)} ausgelegt waren. In der Nähe
half ein \VUgras{Verkäufer} \textsuperscript{(61)} einer \VUgras{Kundin}
\textsuperscript{(62)}, feine seidene \VUgras{Stoffe} \textsuperscript{(64)}
auszusuchen.

%%K t16-12-2 p
Als ich die Treppe hinaufstieg, fiel mir auf, dass der Laden mit japanischen
\VUgras{Lampions} \textsuperscript{(65)}, \VUgras{Sonnenschirmen}
\textsuperscript{(66)} und riesigen \VUgras{Palmen} \textsuperscript{(70)} geschmückt
war.

%%K t16-13-1 p
13. --- Im ersten Stock gab es nicht viel zu sehen: nur Möbel (*),
\VUgras{Geldschränke} \textsuperscript{(67)}, \VUgras{Kommoden}
\textsuperscript{(68)} und \VUgras{Kinderwagen} \textsuperscript{(69)}. Aber der
Gesamtblick über den Laden war sehr \VUgras{unterhaltsam}.

%%K t16-14-1 p
14. --- Unter uns konnte ich die Musikinstrumente sehen: \VUgras{Harfen}
\textsuperscript{(71)}, \VUgras{Gitarren} \textsuperscript{(72)},
\VUgras{Mandolinen} \textsuperscript{(73)}, \VUgras{Kornette}
\textsuperscript{(74)}, \VUgras{Klarinetten} \textsuperscript{(75)}, Geigen mit
ihren \VUgras{Bögen} \textsuperscript{(76)} und sogar eine \VUgras{große Trommel}
\textsuperscript{(77)}. Etwas weiter, am Schreibwarentisch, feilschte ein
\VUgras{Student} \textsuperscript{(78)} mit dem \VUgras{Verkäufer}
\textsuperscript{(79)} um eine Aktentasche. Neben ihnen konnte ich eine schöne
\VUgras{Fibel} \textsuperscript{(80)} erkennen.

%%K t16-14-2 p
Mitten im Laden stand ein hoher \VUgras{Weihnachtsbaum} \textsuperscript{(81)}, über
und über behängt mit Kerzen, Schmuck und Spielzeug jeder Art:
\VUgras{Schaukelpferden} \textsuperscript{(82)}, \VUgras{Puppen}
\textsuperscript{(83)} und so weiter.

%%K t16-14-3 p
Die \VUgras{Parfümerie} \textsuperscript{(84)} mit ihren \VUgras{Mundwässern} und
ihren verschiedenen \VUgras{Parfüms} verbreitete einen sehr angenehmen Duft, der zu
uns heraufstieg.

%%K t16-tit-5 sub
\VUsaut{3.4mm}
\VUpk{10.2pt}{40}{Wir kaufen etwas.}
\VUsaut{2.8mm}

%%K t16-15-1 p
15. --- Meinem Großvater wurde die Zeit lang, und so gingen wir die große Treppe
wieder hinunter. Er blieb einen Augenblick vor einer \VUgras{Sternkarte}
\textsuperscript{(85)} stehen, die, wie er mir sagte, \VUgras{Kometen}
\textsuperscript{(a)}, \VUgras{Mond- und Sonnenfinsternisse} \textsuperscript{(b)},
eine Windrose mit den \VUgras{vier Himmelsrichtungen} \textsuperscript{(c)}
\VUgras{(Norden, Süden, Osten und Westen)}, eine Karte der beiden
\VUgras{Halbkugeln} \textsuperscript{(d)}, die \VUgras{Planeten}
\textsuperscript{(e)} und das \VUgras{Sonnensystem} \textsuperscript{(f)} zeigte. Ich
habe nichts davon verstanden, und mich interessierten viel mehr die betresste Livree
eines \VUgras{Laufburschen} \textsuperscript{(86)}, der vorbeiging, und die hübschen
\VUgras{Fächer} \textsuperscript{(87)} neben mir, bei den \VUgras{Geldbörsen}
\textsuperscript{(88)} und den \VUgras{Brieftaschen} \textsuperscript{(89)}.

%%K t16-16-1 p
16. --- Ich bewunderte auf dem Ladentisch auch ein paar \VUgras{Federn}
\textsuperscript{(90)} und \VUgras{Gürtelschnallen} \textsuperscript{(91)} und die
hübschen \VUgras{künstlichen Blumen} \textsuperscript{(94)}, die anmutig in Körben
angeordnet waren. Die \VUgras{Veilchen}, \VUgras{Goldlack},
\VUgras{Hyazinthen} \textsuperscript{(95)}, \VUgras{Geranien}
\textsuperscript{(96)}, \VUgras{Lilien} \textsuperscript{(97)} und
\VUgras{Kapuzinerkressen} \textsuperscript{(98)} waren sehr gut nachgemacht.

%%K t16-tit-6 sub
\VUsaut{4.4mm}
\VUpk{10.2pt}{40}{Das Geschenk des Großvaters.}
\VUsaut{2.8mm}

%%K t16-17-1 p
17. --- Ich bat meinen Großvater, mir eine der \VUgras{Zahnbürsten}
\textsuperscript{(92)} aus einer Schachtel neben den \VUgras{Haarnadeln}
\textsuperscript{(93)} zu kaufen. Er war einverstanden, und ich ging selbst an die
Kasse, um meinen Einkauf zu bezahlen, neben der gerade ein großes \VUgras{Paket}
\textsuperscript{(100)} für eine \VUgras{Kundin} \textsuperscript{(99)}
zusammengestellt wurde.

%%K t16-18-1 p
18. --- Plötzlich entstand ein kleiner Auflauf unter den Leuten, die bei den
künstlerischen \VUgras{Bronzen} \textsuperscript{(101)} standen. Ein \VUgras{Dieb}
\textsuperscript{(102)} war eben von einem \VUgras{Ladendetektiv}
\textsuperscript{(103)} auf frischer Tat ertappt worden, als er jemanden bestehlen
wollte. Ein Polizist wurde geholt, um ihn zu \VUgras{verhaften}, und gerade als wir
gingen, wurde der Täter zur Wache abgeführt.
\VUsaut{6.0mm}
\VUfilet{22mm}
